\section{有限轴向见证的刚性与八点构型}\label{wr:section}

有限见证必须保留独立尺度。下面先从四轴容量导出任意输入的近矩形结构，
再计算一类八点的矩形阈值，最后用角区构造说明近矩形误差为何不能直接舍去。
本节始终取 $n=2$、$\beta=1$、中心轴平行正方形基。

\begin{lemma}[四轴的精确层容量]\label{wr:layers}
设四中心为 $(\pm t,0),(0,\pm t)$，横轴两个见证边长均为
$X=1+\lambda_x t$，纵轴两个边长均为 $Y=1+\lambda_y t$，
$|\lambda_x|,|\lambda_y|\le1$、$0<t<1/10$。
令 $N$ 为四个正方形示性函数之和，$V_j=|\{N\ge j\}|$，则
\[
 V_2=XY,\quad V_3=X^2+Y^2-XY-4t^2,\quad V_4=(X-2t)(Y-2t).
\]
且对任意 $|E|=1$，记 $I=\int_E N$，有
\begin{align*}
 I&\le3+V_4,\\
 I&\le1+V_2+V_3+V_4.
\end{align*}
两个上界与严格阈值和 $2(1-t)(X+Y)$ 的精确差依次为
\[
 -(\lambda_x+\lambda_y)t+(\lambda_x\lambda_y+4)t^2,
 \qquad
 (\lambda_x+\lambda_y)t+
 (\lambda_x^2+\lambda_y^2+\lambda_x\lambda_y)t^2.
\]
\end{lemma}
\begin{proof}
因 $|X-Y|\le2t$，四个相邻轴点对的交集之并恰为
$R=[-Y/2,Y/2]\times[-X/2,X/2]$；两个相反点对的交集也包含于 $R$。
故 $\{N\ge2\}=R$。
包含横轴点对的两个三重交集之并是边长 $(X-2t),X$ 的竖矩形；
包含纵轴点对的对应并是边长 $Y,(Y-2t)$ 的横矩形。
二者交集正是四重交集，边长为 $(X-2t),(Y-2t)$，从而得到三个面积公式。

逐点有 $N\le3+\mathbf1_{\{N=4\}}$ 及
$N\le1+\sum_{j=2}^4\mathbf1_{\{N\ge j\}}$。
在 $E$ 上积分，利用 $|E|=1$ 即得两条预算；展开后得到最后两式。
\end{proof}

\begin{theorem}[四轴见证迫使二阶近矩形]\label{wr:rigidity}
固定 $M>0$。存在 $C_M,t_M>0$，使以下结论成立。
设 $0<t<t_M$、$\alpha=1-t$、$|E|=1$，并且
\[
 (\pm A_x,0),(0,\pm A_y)\in H_{1,\mathcal Q_c}(\alpha,E),\qquad
 |A_x-t|+|A_y-t|\le Mt^2.
\]
各点任取严格见证边长 $\ell_{\pm x},\ell_{\pm y}$。
则存在 $\rho_t\in[-1,1]$，使
\[
 |E\mathbin\triangle R_{\rho_t,t}|\le C_Mt^2,\qquad
 R_{\rho,t}=[-e^{\rho t}/2,e^{\rho t}/2]\times[-e^{-\rho t}/2,e^{-\rho t}/2],
\]
并且
\[
 \ell_{\pm x}=1-\rho_t t+O_M(t^2),\qquad
 \ell_{\pm y}=1+\rho_t t+O_M(t^2).
\]
输入允许无界；常数统一于所有满足假设的输入和见证选择。
\end{theorem}
\begin{proof}
每个边长都在 $(\alpha,\alpha^{-1})$ 内。
对横轴相反两点，记 $z=(\ell_{+x}+\ell_{-x})/2$、
$u=2A_x$、$d=|\ell_{+x}-\ell_{-x}|$。
因 $d/2\le t+O(t^2)<u$，其精确交叠面积为 $(z-u)(z-d/2)$。
严格两点预算给出
\[
 (z-u)d/2<1+(z-u)z-2\alpha z=O_M(t^2).
\]
这里 $z-u\ge1/2$，最后的统一估计由
$1+(z-u)z-2\alpha z=(z-1)^2-(u-2t)z$ 得到。
故 $\ell_{+x}-\ell_{-x}=O_M(t^2)$；纵轴同理。

把各对尺度替换为平均值，再将 $(\ell-1)/t$ 裁剪到 $[-1,1]$，
并把中心 $A_x,A_y$ 替换为 $t$，每个正方形的对称差面积仅为 $O_M(t^2)$。
因此得到引理~\ref{wr:layers}的四个标准正方形，使每个捕获量仍至少为
$(1-t)\ell-O_M(t^2)$。
记 $L=\lambda_x+\lambda_y$。
两条精确预算分别给 $L\le C_Mt$ 和 $L\ge-C_Mt$，故 $L=O_M(t)$。

还要把预算传回输入。令 $D_3=3+V_4-I$、
$D_1=1+V_2+V_3+V_4-I$；上述估计使 $0\le D_1,D_3\le C_Mt^2$。
精确分解为
\[
 D_3=\int_E(3-N+\mathbf1_{\{N=4\}})+|\{N=4\}\setminus E|,
\]
\[
 D_1=|E\cap\{N=0\}|+\int_{E^c}(N-1)_+.
\]
于是 $|E\setminus R|\le D_3/2$、$|R\setminus E|\le D_1$，其中
$R=\{N\ge2\}=[-Y/2,Y/2]\times[-X/2,X/2]$。
故 $|E\triangle R|=O_M(t^2)$。
又 $XY=1+O_M(t^2)$，将 $R$ 等比例归一化到面积一仅改变 $O_M(t^2)$ 测度。
其短、长边的一阶系数互为相反数；写为 $e^{\pm\rho_t t}$，
再将 $\rho_t$ 裁剪至 $[-1,1]$，仍只改变 $O_M(t^2)$。
由 $Y=1+\lambda_y t$、$X=1+\lambda_x t$ 得到所述尺度展开。
\end{proof}

\begin{corollary}[对角尺度的必要窗口]\label{wr:diagonal-window}
在前一定理假设下，若某个 $(\pm B,\pm B)$ 也是严格光晕点，
$|B-t/2|\le Mt^2$，其任一见证边长写为 $\ell=1+\lambda t$，则
\[
 -1+|\rho_t|-C_Mt\le\lambda\le1-|\rho_t|+C_Mt.
\]
四个对角见证可有不同尺度。
\end{corollary}
\begin{proof}
将 $E$ 换成 $R_{\rho_t,t}$ 的捕获误差至多 $C_Mt^2$。
一维交长公式对所有 $\rho_t,\lambda$ 一致给出
\[
 \frac{|R_{\rho_t,t}\cap Q_\ell(B,B)|}{\ell}
 =1-t\left[\max\{1/2,|\lambda-\rho_t|/2\}
 +\max\{1/2,|\lambda+\rho_t|/2\}\right]+O_M(t^2).
\]
严格阈值使方括号至多为 $1+C_Mt$。
两项各至少 $1/2$，故 $|\lambda\pm\rho_t|\le1+C_Mt$，等价于结论。
\end{proof}

\begin{proposition}[近矩形误差阶不可提高]\label{wr:sharpness}
定理~\ref{wr:rigidity}中的 $O_M(t^2)$ 一般不能改为 $o(t^2)$，
即使再要求四个 $(\pm t/2,\pm t/2)$ 同时属于光晕。
\end{proposition}
\begin{proof}
取 $A_x=A_y=t-t^2$、$Q_0=[-1/2,1/2]^2$，从其中心移走矩形
$D_t=[-t/2,t/2]\times[-t/4,t/4]$，再把该矩形平移到 $(10,10)$ 附近，得到 $E_t$。
输入仍有面积一。四个轴点用边长一见证，捕获量精确为
$1-(t-t^2)-t^2/2=1-t+t^2/2>\alpha$。
四个对角点用边长 $1+t$ 见证，捕获 $1-t^2/2>\alpha(1+t)=1-t^2$。
所有 $R_{\rho,t}$（$|\rho|\le1$）都不包含远置的 $t^2/2$ 质量，
且与 $E_t$ 同为面积一，故 $|E_t\triangle R_{\rho,t}|\ge t^2$。
\end{proof}

\subsection{轴点与对角点的矩形精确阈值}

\begin{theorem}[八点矩形模型的完整分支]\label{wr:eight-rectangle}
令 $1/\sqrt2\le\alpha<1$、$c=2\alpha^2-1$，并置
\[
 A_0=(\alpha^{-1}-\alpha)/2,\quad B_0=(\alpha-\sqrt c)/2,
\]
\[
 \mathscr A_h=\{(\pm hA_0,0),(0,\pm hA_0),(\pm hB_0,\pm hB_0)\},\quad0\le h\le1.
\]
记 $\alpha_c\in(3/4,4/5)$ 为 $4\alpha_c^3-\alpha_c-1=0$ 的唯一根，
$w_*=2\alpha^2/(\alpha+\sqrt c)$、$\widehat w=\min(1,w_*)$。
存在面积一轴平行矩形，使 $\mathscr A_h$ 全部属于其严格光晕，当且仅当
\[
 h<h_*(\alpha),\qquad
 h_*(\alpha)=\frac{(\widehat w+\widehat w^{-1})/2-\alpha}{A_0}
 \wedge\frac{1+\widehat w^2-2\alpha\widehat w}{4\alpha B_0}.
\]
允许输入任意平移和交换两坐标也不改变结论。
若 $h=1-\sigma(1-\alpha)$，则临界收缩量为
\[
 \Sigma_{\rm rect}(\alpha)=
 \begin{cases}
 \dfrac{\alpha^2-\sqrt c}{\alpha(1-\alpha^2)},&1/\sqrt2\le\alpha\le\alpha_c,\\[2mm]
 \dfrac{(w_*-\alpha)(1-\alpha w_*)}{w_*(1-\alpha^2)(1-\alpha)},&\alpha_c\le\alpha<1,
 \end{cases}
\]
严格实现条件是 $\sigma>\Sigma_{\rm rect}(\alpha)$；临界本身不可取。
\end{theorem}
\begin{proof}
归一化矩形为 $R_w$，边长 $w^{-1},w$，$\alpha<w\le1$。
当 $w\le\alpha$ 时光晕为空。
因为八点的每个轨道包含全部坐标反射，且矩形交叠关于中心的每个绝对坐标单调不增，
任何平移矩形若能同时见证八点，其同形中心矩形也能见证。
故可固定矩形中心在原点。

由定理~\ref{rect:exact}证明中的接触尺度分析，
轴向及正对角方向的精确阈值半径分别为
\[
 a_A(w)=\frac{w+w^{-1}}2-\alpha=\frac{N(w)}{2w},\qquad
 b_D(w)=\frac{N(w)}{4\alpha},\quad N(w)=1+w^2-2\alpha w.
\]
第一式的两个轴方向分别使用尺度 $w$ 与 $w^{-1}$。
对角阈值在尺度 $w+2b_D$ 处达到，因为
\[
 b_D-\frac{w^{-1}-w}{4}=\frac{(w-\alpha)(1+w^2)}{4\alpha w}>0.
\]
在另一个内接触 $w^{-1}-2b_D$ 处，平均为
$\alpha-a_A(w)(w-\alpha)/\alpha<\alpha$；
双部分交叠面板上的凸商及其他面板的单调性排除了更大值。
这也直接验证临界 $\alpha=1/\sqrt2$：$w>\alpha$ 仍保证 $w^{-1}/2<\alpha$，
无需交换输入极限和严格水平集。

因此闭半径预算为 $\min(a_A/A_0,b_D/B_0)$，严格见证须严格小于该预算。
$a_A$ 随 $w$ 递减，$b_D$ 随 $w$ 递增，且 $b_D/a_A=w/(2\alpha)$。
唯一交点为 $w_*=2\alpha B_0/A_0=2\alpha^2/(\alpha+\sqrt c)>\alpha$。
若交点超过一，最大值位于 $w=1$；否则位于交点。
条件 $w_*\le1$ 等价于 $4\alpha^3-\alpha-1\ge0$，得到两分支。
将 $\Sigma=(1-h_*)/(1-\alpha)$ 化简即得公式。
\end{proof}

\begin{proposition}[阈值的任意阶展开及显式余项]\label{wr:eight-series}
高水平分支在 $t=1-\alpha=0$ 附近解析，且
\[
 \Sigma_{\rm rect}(1-t)=\frac38+\frac58t+\frac{29}{32}t^2+
 \frac{17}{16}t^3+\frac{51}{128}t^4-\frac{471}{128}t^5+O(t^6).
\]
令 $w(t)=\sum_{j\ge0}w_jt^j$、$w_0=1$，则其系数由
\[
 (2t-t^2)w^2-4(1-t)^3w+4(1-t)^4=0
\]
唯一递推得到。具体地，对 $n\ge1$、$W_{n-1}=\sum_{j<n}w_jt^j$，
\[
 w_n=\frac14[t^n]\left\{(2t-t^2)W_{n-1}^2
 -4[(1-t)^3-1]W_{n-1}+4(1-t)^4\right\}.
\]
把 $w$ 代回前一定理得到任意阶系数 $s_j$。
对任意 $N\ge0$、$0\le t<1/8$，有显式尾界
\[
 \left|\Sigma_{\rm rect}(1-t)-\sum_{j=0}^Ns_jt^j\right|
 \le\frac{4(8t)^{N+1}}{1-8t}.
\]
例如常数截断和线性截断的误差阶分别恰为 $t$ 和 $t^2$。
\end{proposition}
\begin{proof}
隐式方程对 $w$ 的导数在 $(t,w)=(0,1)$ 为 $-4$，故存在唯一解析分支。
配系数得到递推式和所列展开。
为给出可复核余项，令 $a=1-t$，在复圆 $|t|=1/8$ 上有
$|2a^2-2|\le17/32$，因此主平方根满足
$|a+\sqrt{2a^2-1}|>3/2$，且 $|w|>1/2$。
再置 $u=a(1+a)/(a+\sqrt{2a^2-1})^2$，则 $|u|\le17/16$，
\[
 \Sigma=\frac{u(1+a-au)}{w(1+a)},\qquad |\Sigma|<4.
\]
该表达式也消去了 $t=0$ 的可去奇点。
Cauchy 系数界给 $|s_j|\le4\,8^j$，求几何尾和即得结论。
首两个遗漏系数 $5/8$、$29/32$ 非零，故所述误差阶不能提高。
\end{proof}

\subsection{角区质量不能从八点见证中舍去}

\begin{theorem}[非矩形输入改善另一八点族]\label{wr:l-counterexample}
令
\[
 q=\frac{4-\sqrt7}{6},\quad p=2q-q^2,\quad
 \gamma=2q-2q^2=\frac{1+2\sqrt7}{18}<\frac38.
\]
固定 $\gamma<\sigma<3/8$、$\alpha=1-t$，并置
\[
 r_0=\frac{1-\alpha^2}{\alpha+\sqrt{\alpha^2-\tfrac34(1-\alpha^2)}},\qquad
 r=r_0(1-\sigma t),\qquad
 \mathscr X_r=\{(\pm3r/4,\pm r/4),(\pm r/4,\pm3r/4)\}.
\]
当 $t>0$ 充分小时，存在面积一的有界轴平行多边形输入 $E_t$，
使 $\mathscr X_r$ 全部属于其严格光晕；但不存在任何面积一轴平行矩形输入实现同一八点族。
$t$ 的阈值在 $\sigma$ 的任意紧子区间上一致。
\end{theorem}
\begin{proof}
定义
\[
 \ell=\sqrt{1+(1/4+4p)r^2}-r,
 \qquad R=[-(\ell+3r/2)/2,(\ell+3r/2)/2]
 \times[-(\ell+r/2)/2,(\ell+r/2)/2].
\]
在 $R$ 每个角用朝内坐标 $(u,v)$ 移去 L 形区域
\[
 0\le u\le3r/2,\quad0\le v\le r/2,\quad
 (u\le qr\ \text{或}\ v\le qr).
\]
每孔面积为 $pr^2$，四孔在小 $t$ 下互不相交，且
\[
 |E_t|=(\ell+3r/2)(\ell+r/2)-4pr^2=1.
\]
对 $(\pm3r/4,\pm r/4)$ 取尺度 $\ell$。
每个见证正方形只遇到对应角的孔，捕获量为 $\ell^2-pr^2$。
对另外四点取尺度 $\ell_2=\ell+2(1-q)r$。
其与 $R$ 的交叠为边长 $\ell+(3/2-q)r$、$\ell+(1/2-q)r$ 的矩形。
四孔被捕获的总面积系数为
\[
 p+(3q/2-q^2)+(q/2-q^2)=4q-3q^2=3/4.
\]
因此两组见证的严格捕获余量分别是
\[
 \Delta_1=\ell(\ell-\alpha)-pr^2,\qquad
 \Delta_2=(\ell-\alpha)\ell_2-pr^2.
\]
由 $r_0=t+7t^2/8+O(t^3)$，直接展开得
$\Delta_1=(\sigma-\gamma)t^2+O_\sigma(t^3)>0$。
这还强制 $\ell>\alpha$，从而 $\Delta_2>\Delta_1>0$。
故同一输入严格见证八点，且所有边长自动落在合法窗口。

最后排除矩形。平移不能改善完整反射轨道的共同可行性，故仍只需考察中心矩形 $R_w$、
$\alpha<w\le1$。令 $A=(w+w^{-1})/2-\alpha$。
由定理~\ref{rect:exact}中的象限边界，点 $(3r/4,r/4)$ 满足 $y-x\le(w^{-1}-w)/2$，
故属于光晕必须有
\[
 r<\frac{1+w^2-2\alpha w}{w+\alpha}=:r_1(w).
\]
而
\[
 r_1(w)-A_0=
 \frac{(w-\alpha)(2\alpha w-\alpha^2-1)}{2\alpha(w+\alpha)}<0.
\]
坐标交换不改变八点族。
另一方面 $r-A_0=(3/8-\sigma)t^2+O_\sigma(t^3)>0$，故任何这样的矩形都不可能实现八点。
此反例只比较有限点的共同见证，不给出全体光晕体积的三阶下界。
\end{proof}
